% --- Archon physics formalization source begin ---
% archon:physics
% archon:covers IPhO2026Problems/problem_IPhO_2026_2_C_2.lean
% archon:source-report reports/ipho_2026_k3/problem_IPhO_2026_2_C_2.source.json
% archon:problem-id IPhO_2026_2
% archon:part-id C.2
% archon:previous-part IPhO_2026_2_C_1
% archon:previous-part-policy natural_language_prerequisite_only

\chapter{Physics problem IPhO\_2026\_2\_C\_2}
\label{ch:IPhO2026Problems_problem_IPhO_2026_2_C_2}

\paragraph{Problem source.}
For the half-cylindrical mirror of radius R, ray A is incident at angle theta
and its reflected line is y = m\_A*x + b\_A.  A neighboring parallel ray B is
incident at theta + Delta theta, with Delta theta much smaller than theta, and
its reflected line is y = m\_B*x + b\_B.  The envelope/intersection of neighboring
rays forms the caustic.  Use Figure 2g and its coordinate convention.

Current subquestion:
Expand m\_B and b\_B to first order in Delta theta.

\paragraph{Current subquestion.}
Expand m\_B and b\_B to first order in Delta theta.

\paragraph{Recorded answer/context.}
m\_B = cot(2*theta) - 2*csc(2*theta)\textasciicircum{}2*Delta theta; b\_B = [R/(2*cos(theta))]*(1 + tan(theta)*Delta theta), up to O(Delta theta\textasciicircum{}2).

\paragraph{Figure/image path.}
/root/proposal\_for\_physic/science-mango/ipho\_2026\_source/image/T2\_page-4.png

\paragraph{Reusable previous-part conclusions.}
\begin{itemize}
\item Source C.1. Question: Write the slope m\_A and intercept b\_A of reflected ray A in terms of theta and R. Reusable conclusions: m\_A = cot(2*theta), and b\_A = R/(2*cos(theta)). Policy: natural\_language\_prerequisite\_only; do\_not\_import\_Lean\_output
\end{itemize}

\paragraph{Formalization target.}
create a compiling Lean file with sorry bodies at `IPhO2026Problems/problem\_IPhO\_2026\_2\_C\_2.lean`.
The Lean declarations must preserve the physical quantities, dimensions or dimensional roles, figure labels, governing-law hypotheses, and final relation expressed by this problem.
Use Mathlib/Physlib names found through LeanExplore where available. If a domain API is missing, introduce faithful local abstractions rather than scalar placeholder aliases.

\begin{theorem}[Physics formalization target]
\label{thm:physics:IPhO_2026_2_C_2:target}
\uses{thm:IPhO2026Problems_problem_IPhO_2026_2_C_2:ray_B_first_order_expansion}
The assigned autoformalize agent should translate this physics problem into Lean declarations in the covered file, with theorem and lemma proof bodies written as `by sorry`.
\end{theorem}
\begin{proof}
This is an autoformalization task, not a proof task. Produce faithful statements that can later be proved without weakening the source contract.
\end{proof}
% NOTE: PhysLean-coverage exemption (planner-recorded, iter-002): PhysLean has no geometric-optics/asymptotic-expansion module for the first-order mirror regime; the formalization uses Mathlib `Asymptotics.IsLittleO` contracts (see the physics-grounding log for this file). Self-containment is kept with the `import Mathlib` baseline; no irrelevant Physlib import is added. This documents the resolved import policy (iter-001 review ruling) for the blueprint-doctor `missing-physlib-import` check.
% --- Archon physics formalization source end ---

% --- Archon named-quantities coverage (blueprint-writer 2-c-2-entries) ---
\subsection*{Governing-law structure: the neighboring-ray expansion package}

\begin{definition}[Neighboring-ray expansion package]
\label{def:IPhO2026Problems_problem_IPhO_2026_2_C_2:NeighboringRayExpansion}
\lean{IPhO2026_2_C_2.NeighboringRayExpansion}
The Figure-2g C.2 package: the half-cylindrical mirror of radius $R > 0$;
the one-parameter reflected slope--intercept family
$\varphi \mapsto (M\,\varphi, B\,\varphi)$ generated by specular reflection
of the axial (parallel-to-$y$) rays, with ray A the member at
$\varphi = \theta$ and the neighboring parallel ray B the member at
$\varphi = \theta + \Delta\theta$; the acute branch datum
$0 < \theta < \pi/2$ keeping the mirror hit on the upper right quarter; the
C.1 recorded values of $m_A(\theta)$ and $b_A(\theta)$ as natural-language
prerequisite hypotheses; and the smooth-governing-law regularity interface,
certifying that $M$ and $B$ each admit a first-order expansion about
$\theta$ with remainder little-$o$ of $\Delta\theta$ (a Mathlib
\texttt{Asymptotics.IsLittleO} contract over the neighborhood filter of
$0$).  All coordinate readouts live in the Figure-2g plane: $R$ and the
intercepts carry the dimension of length, the slopes and angles are
dimensionless.
\end{definition}
\begin{proof}
Definition; no claim.
\end{proof}

\subsection*{Derivation bridge: the Figure-2g branch is nonsingular}

\begin{lemma}[Branch denominators are nonzero]
\label{lem:IPhO2026Problems_problem_IPhO_2026_2_C_2:branch_denominators_ne_zero}
\lean{IPhO2026_2_C_2.NeighboringRayExpansion.branch_denominators_ne_zero}
\uses{def:IPhO2026Problems_problem_IPhO_2026_2_C_2:NeighboringRayExpansion}
On the acute Figure-2g branch of the expansion package, the denominators
appearing in the C.1/C.2 formulas do not vanish:
$\sin (2\theta) \neq 0$ and $\cos \theta \neq 0$, so $\cot(2\theta)$,
$\csc(2\theta)$, $\tan\theta$ and the C.1 intercept are all well-formed
quantities of the subquestion.  Bridge lemma only; not a target conclusion.
\end{lemma}
\begin{proof}
The branch datum gives $0 < \theta < \pi/2$, hence $0 < 2\theta < \pi$;
there $\sin(2\theta) > 0$ and $\cos\theta > 0$, so both are nonzero.
\end{proof}

\subsection*{Value theorems: the C.2 first-order expansions}

\begin{theorem}[Neighboring-ray slope to first order]
\label{thm:IPhO2026Problems_problem_IPhO_2026_2_C_2:ray_B_slope_first_order}
\lean{IPhO2026_2_C_2.NeighboringRayExpansion.ray_B_slope_first_order}
\uses{def:IPhO2026Problems_problem_IPhO_2026_2_C_2:NeighboringRayExpansion, lem:IPhO2026Problems_problem_IPhO_2026_2_C_2:branch_denominators_ne_zero, lem:IPhO2026Problems_problem_IPhO_2026_2_C_2:slope_deriv_value, def:IPhO2026Problems_problem_IPhO_2026_2_C_2:slopeFamily_deriv}
\textbf{(Target, T2-C.2, slope half; coefficients conclusion-side.)}  As
$\Delta\theta \to 0$,
\[
m_B(\theta, \Delta\theta)
= \cot (2\theta) - 2\,\csc(2\theta)^{2}\,\Delta\theta + o(\Delta\theta) ,
\]
stated formally as an \texttt{IsLittleO} contract of the defect
$\Delta\theta \mapsto m_B(\theta,\Delta\theta) -
(\cot(2\theta) - 2\,\csc(2\theta)^2\,\Delta\theta)$ against
$\Delta\theta \mapsto \Delta\theta$ in the neighborhood filter of $0$.
\end{theorem}
\begin{proof}
Family membership rewrites $m_B(\theta,\Delta\theta) = M(\theta+\Delta\theta)$
and the smooth-governing-law interface expands
$M(\theta+\Delta\theta) = M(\theta) + (\partial_\theta M)\,\Delta\theta
+ o(\Delta\theta)$.  The C.1 value $M(\theta) = \cot(2\theta)$ fixes the
zeroth order; the derivative identity
$(d/d\theta)\cot(2\theta) = -2\csc(2\theta)^2$ of the specular family,
nonsingular on the branch by
\cref{lem:IPhO2026Problems_problem_IPhO_2026_2_C_2:branch_denominators_ne_zero},
fixes the first-order coefficient.
\end{proof}

\begin{theorem}[Neighboring-ray intercept to first order]
\label{thm:IPhO2026Problems_problem_IPhO_2026_2_C_2:ray_B_intercept_first_order}
\lean{IPhO2026_2_C_2.NeighboringRayExpansion.ray_B_intercept_first_order}
\uses{def:IPhO2026Problems_problem_IPhO_2026_2_C_2:NeighboringRayExpansion, lem:IPhO2026Problems_problem_IPhO_2026_2_C_2:branch_denominators_ne_zero, lem:IPhO2026Problems_problem_IPhO_2026_2_C_2:intercept_deriv_value, def:IPhO2026Problems_problem_IPhO_2026_2_C_2:interceptFamily_deriv}
\textbf{(Target, T2-C.2, intercept half; coefficients conclusion-side.)}
As $\Delta\theta \to 0$,
\[
b_B(\theta, \Delta\theta)
= \frac{R}{2\cos\theta}\,\bigl(1 + \tan\theta\,\Delta\theta\bigr)
+ o(\Delta\theta) ,
\]
stated formally as an \texttt{IsLittleO} contract of the defect against
$\Delta\theta \mapsto \Delta\theta$ in the neighborhood filter of $0$.
\end{theorem}
\begin{proof}
Family membership rewrites $b_B(\theta,\Delta\theta) = B(\theta+\Delta\theta)$
and the interface expands $B$ to first order with little-$o$ remainder.
The C.1 value $B(\theta) = R/(2\cos\theta)$ fixes the zeroth order; the
derivative identity
$(d/d\theta)\bigl(R/(2\cos\theta)\bigr)
= \bigl(R/(2\cos\theta)\bigr)\tan\theta$ of the specular family fixes the
first-order coefficient.
\end{proof}

\begin{theorem}[Neighboring-ray first-order expansion (C.2 main target)]
\label{thm:IPhO2026Problems_problem_IPhO_2026_2_C_2:ray_B_first_order_expansion}
\lean{IPhO2026_2_C_2.NeighboringRayExpansion.ray_B_first_order_expansion}
\uses{thm:IPhO2026Problems_problem_IPhO_2026_2_C_2:ray_B_slope_first_order, thm:IPhO2026Problems_problem_IPhO_2026_2_C_2:ray_B_intercept_first_order}
\textbf{(Target, T2-C.2.)}  The two recorded C.2 expansions hold jointly:
the reflected slope and intercept of the neighboring ray B expand to first
order in $\Delta\theta$ with remainders little-$o$ of $\Delta\theta$ ---
the faithful first-order content of the official $+\,O(\Delta\theta^{2})$
phrasing.  The expansion coefficients appear only here and in the two halves
above, never hypothesis-side.
\end{theorem}
\begin{proof}
Conjunction of
\cref{thm:IPhO2026Problems_problem_IPhO_2026_2_C_2:ray_B_slope_first_order}
and
\cref{thm:IPhO2026Problems_problem_IPhO_2026_2_C_2:ray_B_intercept_first_order};
the formal proof is exact pair introduction of the two halves.
\end{proof}

% --- Iter-012 ledger transcription (plan agent): iter-011 prover-created ---
% support declarations of the HasDerivAt redraft; statements/routes read
% first-hand from the Lean file.  These entries carry the dependency edges
% of the two honesty-discount sorries in the target theorems.

\subsection*{Specular-family evaluation layer (iter-011 redraft helpers)}

\begin{definition}[Specular slope family (C.1 closed form)]
\label{def:IPhO2026Problems_problem_IPhO_2026_2_C_2:specularSlopeFamily}
\lean{IPhO2026_2_C_2.NeighboringRayExpansion.specularSlopeFamily}
\uses{def:IPhO2026Problems_problem_IPhO_2026_2_C_2:NeighboringRayExpansion}
The closed-form reflected slope of the specular ray family of the
half-cylindrical mirror, $\varphi \mapsto \cos(2\varphi)/\sin(2\varphi)$
($= \cot(2\varphi)$) --- the C.1 law of reflection phrased as a quotient
so the Mathlib derivative chain rule applies.  Used to \emph{evaluate}
derivatives; the C.2 expansion coefficients stay conclusion-side.
\end{definition}
\begin{proof}
Definition; no claim.
\end{proof}

\begin{definition}[Specular intercept family (C.1 closed form)]
\label{def:IPhO2026Problems_problem_IPhO_2026_2_C_2:specularInterceptFamily}
\lean{IPhO2026_2_C_2.NeighboringRayExpansion.specularInterceptFamily}
\uses{def:IPhO2026Problems_problem_IPhO_2026_2_C_2:NeighboringRayExpansion}
The closed-form reflected intercept of the specular ray family,
$\varphi \mapsto (R/2)\,(\cos\varphi)^{-1}$ ($= R/(2\cos\varphi)$), the
C.1 intercept law phrased for the chain rule.
\end{definition}
\begin{proof}
Definition; no claim.
\end{proof}

\begin{definition}[Specular-family deriv contract (slope)]
\label{def:IPhO2026Problems_problem_IPhO_2026_2_C_2:slopeFamily_deriv}
\lean{IPhO2026_2_C_2.NeighboringRayExpansion.slopeFamily_deriv}
\uses{def:IPhO2026Problems_problem_IPhO_2026_2_C_2:NeighboringRayExpansion, def:IPhO2026Problems_problem_IPhO_2026_2_C_2:specularSlopeFamily}
Bridge hypothesis (a Prop, not a structure field): the deriv-number of
the abstract family $M$ at $\theta$ equals that of the specular C.1
closed form, $\mathrm{deriv}\,M\,\theta =
\mathrm{deriv}\,\texttt{specularSlopeFamily}\,\theta$ --- the statement
that $M$ is, to first order, the specular family.  The two target
theorems each discharge a local instance of it; that discharge is the
slope-side honesty-discount `sorry` recorded in the task result.
\end{definition}
\begin{proof}
Definition; no claim.
\end{proof}

\begin{definition}[Specular-family deriv contract (intercept)]
\label{def:IPhO2026Problems_problem_IPhO_2026_2_C_2:interceptFamily_deriv}
\lean{IPhO2026_2_C_2.NeighboringRayExpansion.interceptFamily_deriv}
\uses{def:IPhO2026Problems_problem_IPhO_2026_2_C_2:NeighboringRayExpansion, def:IPhO2026Problems_problem_IPhO_2026_2_C_2:specularInterceptFamily}
The intercept-side counterpart:
$\mathrm{deriv}\,B\,\theta =
\mathrm{deriv}\,(\texttt{specularInterceptFamily}\,R)\,\theta$;
the second honesty-discount `sorry` of the file discharges a local
instance of this contract.
\end{definition}
\begin{proof}
Definition; no claim.
\end{proof}

\begin{definition}[Specular-family smoothness interface (slope)]
\label{def:IPhO2026Problems_problem_IPhO_2026_2_C_2:slopeFamily_smooth}
\lean{IPhO2026_2_C_2.NeighboringRayExpansion.slopeFamily_smooth}
\uses{def:IPhO2026Problems_problem_IPhO_2026_2_C_2:NeighboringRayExpansion}
Physical smoothness of the specular family on the acute branch:
$M$ is differentiable at every $\varphi \in (0, \pi/2)$ --- the law of
reflection depends smoothly on the incidence point and the branch carries
no singularity ($\cos\varphi \ne 0$, $\sin 2\varphi \ne 0$).
\end{definition}
\begin{proof}
Definition; no claim.
\end{proof}

\begin{definition}[Specular-family smoothness interface (intercept)]
\label{def:IPhO2026Problems_problem_IPhO_2026_2_C_2:interceptFamily_smooth}
\lean{IPhO2026_2_C_2.NeighboringRayExpansion.interceptFamily_smooth}
\uses{def:IPhO2026Problems_problem_IPhO_2026_2_C_2:NeighboringRayExpansion}
The intercept-side counterpart: $B$ is differentiable at every acute
incidence angle.
\end{definition}
\begin{proof}
Definition; no claim.
\end{proof}

\begin{lemma}[Derivative of the specular slope closed form]
\label{lem:IPhO2026Problems_problem_IPhO_2026_2_C_2:deriv_specularSlopeFamily}
\lean{IPhO2026_2_C_2.NeighboringRayExpansion.deriv_specularSlopeFamily}
\uses{def:IPhO2026Problems_problem_IPhO_2026_2_C_2:specularSlopeFamily}
$(d/d\varphi)\,\cot(2\varphi) = -2\csc(2\varphi)^{2}$ at every acute
$\varphi$, nonsingular on the branch.  Proved `sorry`-free in Lean by the
Mathlib chain rule on the quotient form.
\end{lemma}
\begin{proof}
Write $\cot(2\varphi) = \cos(2\varphi)/\sin(2\varphi)$, differentiate by
the quotient rule with $\sin(2\varphi) \ne 0$ on the branch, and simplify
via $\sin^{2} + \cos^{2} = 1$.
\end{proof}

\begin{lemma}[Derivative of the specular intercept closed form]
\label{lem:IPhO2026Problems_problem_IPhO_2026_2_C_2:deriv_specularInterceptFamily}
\lean{IPhO2026_2_C_2.NeighboringRayExpansion.deriv_specularInterceptFamily}
\uses{def:IPhO2026Problems_problem_IPhO_2026_2_C_2:specularInterceptFamily}
$(d/d\varphi)\,(R/(2\cos\varphi)) = (R/(2\cos\varphi))\tan\varphi$ at
every acute $\varphi$.  Proved `sorry`-free in Lean
(chain rule with $\cos\varphi \ne 0$).
\end{lemma}
\begin{proof}
Differentiate $(R/2)(\cos\varphi)^{-1}$ by the reciprocal rule and
rewrite $\sin/\cos = \tan$.
\end{proof}

\begin{lemma}[Slope coefficient of a contracted family]
\label{lem:IPhO2026Problems_problem_IPhO_2026_2_C_2:slope_deriv_value}
\lean{IPhO2026_2_C_2.NeighboringRayExpansion.slope_deriv_value}
\uses{def:IPhO2026Problems_problem_IPhO_2026_2_C_2:NeighboringRayExpansion, def:IPhO2026Problems_problem_IPhO_2026_2_C_2:slopeFamily_deriv, lem:IPhO2026Problems_problem_IPhO_2026_2_C_2:deriv_specularSlopeFamily}
The coefficient $dm$ of any instance of the regularity interface
($\texttt{HasDerivAt}\,M\,dm\,\theta$) whose family satisfies the
specular deriv contract is exactly the recorded C.2 slope coefficient
$-2\csc(2\theta)^{2}$.  `Sorry`-free in Lean.
\end{lemma}
\begin{proof}
$dm = \mathrm{deriv}\,M\,\theta$ by uniqueness of the deriv-number
(\texttt{HasDerivAt.deriv}); rewrite by the contract and apply
\cref{lem:IPhO2026Problems_problem_IPhO_2026_2_C_2:deriv_specularSlopeFamily}.
\end{proof}

\begin{lemma}[Intercept coefficient of a contracted family]
\label{lem:IPhO2026Problems_problem_IPhO_2026_2_C_2:intercept_deriv_value}
\lean{IPhO2026_2_C_2.NeighboringRayExpansion.intercept_deriv_value}
\uses{def:IPhO2026Problems_problem_IPhO_2026_2_C_2:NeighboringRayExpansion, def:IPhO2026Problems_problem_IPhO_2026_2_C_2:interceptFamily_deriv, lem:IPhO2026Problems_problem_IPhO_2026_2_C_2:deriv_specularInterceptFamily}
The coefficient $db$ of any $\texttt{HasDerivAt}\,B\,db\,\theta$ instance
on a contracted family is $(R/(2\cos\theta))\tan\theta$.  `Sorry`-free
in Lean.
\end{lemma}
\begin{proof}
As
\cref{lem:IPhO2026Problems_problem_IPhO_2026_2_C_2:slope_deriv_value}
with
\cref{lem:IPhO2026Problems_problem_IPhO_2026_2_C_2:deriv_specularInterceptFamily}.
\end{proof}

\begin{lemma}[Contract evaluation corollary (slope)]
\label{lem:IPhO2026Problems_problem_IPhO_2026_2_C_2:slope_deriv_of_contract}
\lean{IPhO2026_2_C_2.NeighboringRayExpansion.slope_deriv_of_contract}
\uses{def:IPhO2026Problems_problem_IPhO_2026_2_C_2:slopeFamily_deriv, lem:IPhO2026Problems_problem_IPhO_2026_2_C_2:deriv_specularSlopeFamily}
On a family satisfying the slope deriv contract,
$\mathrm{deriv}\,M\,\theta = -2\csc(2\theta)^{2}$ --- the evaluation
corollary (assumption-side, not a C.2 target).  `Sorry`-free in Lean.
\end{lemma}
\begin{proof}
Rewrite the contract and apply
\cref{lem:IPhO2026Problems_problem_IPhO_2026_2_C_2:deriv_specularSlopeFamily}.
\end{proof}

\begin{lemma}[Contract evaluation corollary (intercept)]
\label{lem:IPhO2026Problems_problem_IPhO_2026_2_C_2:intercept_deriv_of_contract}
\lean{IPhO2026_2_C_2.NeighboringRayExpansion.intercept_deriv_of_contract}
\uses{def:IPhO2026Problems_problem_IPhO_2026_2_C_2:interceptFamily_deriv, lem:IPhO2026Problems_problem_IPhO_2026_2_C_2:deriv_specularInterceptFamily}
On a family satisfying the intercept deriv contract,
$\mathrm{deriv}\,B\,\theta = (R/(2\cos\theta))\tan\theta$.
`Sorry`-free in Lean.
\end{lemma}
\begin{proof}
Rewrite the contract and apply
\cref{lem:IPhO2026Problems_problem_IPhO_2026_2_C_2:deriv_specularInterceptFamily}.
\end{proof}
